% ============================================================================
% ICSE Paper -- Approach Section
%
% v7 rewrite -- updated to track s_linker17e (v2.6.2 production: multi-framing
% extraction with union merge + validated coref).  Changes from v6:
%
%   1. Story arc.  The trace-link definition still pins every link to a
%      textual reference, but the primary mechanism is now the
%      \textbf{three-framing ensemble} (\linkerA mention-framing,
%      \linkerB actor-framing, \linkerC entity-framing): three complementary
%      linguistic angles on the same trace-link evidence, run in parallel,
%      all alias-aware, unioned, and validated by a single evidence-bundle
%      pass.  The supplementary linker is \coref, which closes the
%      pronominal gap with citation grounding plus a coref-specific
%      single-pass validation (the v2.6.2 / s17e change; +3.2pp F1, FP
%      43->14 vs s15).
%   2. Architecture replaces "AnchorLinker per-component disambiguation"
%      (gone in s17e) with "union merge + unified validation."
%   3. Style.  Wang's first-author conventions: single-sentence section
%      scope opener, "\approach in a Nutshell" subsection with
%      (i)/(ii)/(iii) walkthrough, "thereby" consequence closers, bolded
%      inline mini-subheadings, declarative first-person plural, "to the
%      best of our knowledge" positioning, bottom-up exposition.
% ============================================================================

% --------------------------------------------------------------------------
% Naming macros.  Edit a single line here to rename the pipeline or any of
% the linkers; all in-text references update automatically.  \xspace
% absorbs the trailing space when followed by punctuation.
%
% \linkerA/\linkerB/\linkerC name the three linguistic framings that
% together form the primary multi-framing linker.  \coref names the
% supplementary pronominal linker.  \multiframing names the three-framing
% ensemble as a unit.
% --------------------------------------------------------------------------
\providecommand{\approach}{AALinker\xspace}
\providecommand{\multiframing}{multi-framing linker\xspace}
\providecommand{\linkerA}{Framing\,A\xspace}      % explicit-mention framing
\providecommand{\linkerB}{Framing\,B\xspace}      % actor-role framing
\providecommand{\linkerC}{Framing\,C\xspace}      % alias-aware entity framing
\providecommand{\coref}{CorefLinker\xspace}       % supplementary, pronoun links

\section{\approach}\label{sec:approach}

This section presents \approach, a training-free pipeline that
recovers trace links from an architecture document to the components
of an architectural model.
Given a document~$D$ with sentences $S = \{s_1, \ldots, s_n\}$ and a
component set $C = \{c_1, \ldots, c_m\}$, \approach outputs a link set
$L \subseteq S \times C$, where each link $(s_i, c_j)$ asserts that
sentence~$s_i$ discusses the architectural responsibilities, behaviour,
or interactions of component~$c_j$.

Trace links are, by definition, anchored on a textual reference to the
linked component: the linked sentence has to mention the component in
some form.
That form is not unique.
The same trace link surfaces, in our corpus, under three complementary
linguistic forms:
the component appears as an \emph{explicit mention} of its name
(``the \texttt{EventDispatcher} routes requests'');
the component appears as the \emph{actor} of an architectural action
(``\texttt{EventDispatcher} routes requests'');
or the component is named indirectly through a project-specific
\emph{entity alias} --- an abbreviation, a synonym, a partial name, or
a functional role (``the dispatcher routes requests'').
A residual fraction of trace links carry no surface name at all and
refer to the component only through a discourse pronoun
(\emph{it}, \emph{this}, \emph{their}, $\ldots$).\todo{cite empirical
distribution: framings together cover $\sim$95\% of gold links;
pronoun-only links account for $<$5\%.}
Two facts complicate this picture.
First, component names are often ordinary technical English
(\texttt{Core}, \texttt{Adapter}, \texttt{Logic}, \texttt{Storage}), so
a surface-level lexical match overshoots into generic uses of those
words.
Second, project-specific aliases are not known in advance, so any
linker that targets indirect reference forms must first discover them
at runtime.

\approach resolves these complications in three stages.
\textbf{Project knowledge discovery} (\S\ref{sec:knowledge}) learns,
at runtime, which component names are ambiguous and which document
terms are aliases for which components.
The \textbf{\multiframing} (\S\ref{sec:multiframing}) recovers the
majority of trace links by running three complementary linguistic
framings --- \linkerA, \linkerB, and \linkerC --- in parallel,
unioning their outputs into one candidate pool, and applying a single
evidence-bundle validation pass that is the sole quality gate for
multi-framing candidates.
The supplementary \textbf{\coref} (\S\ref{sec:coref}) closes the
residual pronominal gap by resolving discourse references and
verifying each resolution with a deterministic antecedent check and
a coref-specific validation pass.
A trivial merge step (\S\ref{sec:merge}) deduplicates the two link
sets.
To the best of our knowledge, \approach is the first trace-link
recovery pipeline that combines runtime architectural-knowledge
discovery, a definition-driven multi-framing primary linker, and an
evidence-grounded pronominal linker for
architecture-document-to-model tracing.

\subsection{\approach in a Nutshell}\label{sec:nutshell}

\approach is organised as a six-phase directed acyclic graph
(Figure~\ref{fig:pipeline}).
As illustrated in Figure~\ref{fig:pipeline}, \approach proceeds in
six steps:
(i)~two independent knowledge analyses --- model understanding and
document understanding --- run in parallel and emit a
\emph{model ambiguity map} and a \emph{document alias table};
(ii)~the three framings of the \multiframing run in parallel, each
receiving the full alias table, and each emits its own candidate
link set;
(iii)~a \emph{union merge} folds the three framing outputs into a
single candidate pool;
(iv)~a \emph{unified evidence-bundle validation} pass --- consisting
of a generic-word precheck and a two-pass approve-or-reject decision
on architectural participation and referential specificity --- is the
sole quality gate for the multi-framing pool;
(v)~\coref runs in parallel with the multi-framing pipeline, resolves
each pronoun-bearing sentence against a $\pm 5$ context window with
a deterministic antecedent check, and then passes the resolutions
through a coref-specific validation;
(vi)~the two validated link sets are concatenated and deduplicated.
The fan-out/fan-in structure is deliberate: every expensive LLM call
whose inputs are independent runs in parallel, thereby keeping the
end-to-end latency of the pipeline within a single round-trip per
project regardless of how many framings or linkers participate.

Two design rules cut across all six phases.
\textbf{Union for open-ended extraction.}
Where the LLM may miss true references because of a single
prompt-specific bias, \approach unions the outputs of independent
prompts and lets a downstream validation pass do the filtering ---
thereby keeping recall above the cost of any one framing.
\textbf{Evidence for validation.}
Where a decision can be checked against the document, \approach
requires the LLM to expose the evidence it used --- matched spans,
mention types, anchor sentences, or antecedent sentence numbers ---
so that a deterministic check, rather than an opaque second LLM call,
decides acceptance.

\begin{table}[t]
\caption{The two \approach linkers and the framings that compose
them.  The \multiframing unions three linguistic framings and applies
a single unified validation pass; \coref runs its own
extraction-plus-validation pipeline for pronouns.}
\label{tab:strategies}
\centering\footnotesize
\begin{tabular}{@{}llll@{}}
\toprule
\textbf{Linker / Framing} & \textbf{Reference form}   & \textbf{Knowledge consumed}    & \textbf{LLM calls} \\
\midrule
\multicolumn{4}{@{}l@{}}{\emph{\multiframing} -- union merge + unified evidence-bundle validation}\\
\quad\linkerA          & Explicit mention          & Alias table                   & 1 batched          \\
\quad\linkerB          & Actor role                & Alias table                   & 1 batched          \\
\quad\linkerC          & Entity alias / role       & Alias table                   & 2 batched          \\
\quad\textit{validation} & ---                     & Ambiguity map, anchor sents.\ & 2 + 1/ambig.\ comp.\ \\
\midrule
\multicolumn{4}{@{}l@{}}{\emph{\coref} -- cited antecedent + coref-specific validation}\\
\quad extraction        & Pronoun, role-referent   & Alias table, $\pm 5$ window   & 1 batched          \\
\quad validation        & ---                       & Component list                & 1 batched          \\
\bottomrule
\end{tabular}
\end{table}

\begin{figure}[t]
\centering
\fbox{\parbox{0.94\linewidth}{
\small
\textbf{\approach pipeline.}
Phase~1 runs model and document understanding in parallel and
produces an ambiguity map and an alias table.
Phase~2 runs \linkerA, \linkerB, and \linkerC in parallel, all
alias-aware.
Phase~3 union-merges their outputs; Phase~4 validates the union
through one evidence-bundle pass.
Phase~5 (parallel with Phases~2--4) runs \coref's extraction and
its coref-specific validation.
Phase~6 deduplicates the two validated link sets.
}}
\caption{\approach organises knowledge discovery, multi-framing
extraction-and-validation, and pronominal linking as a six-phase
DAG.}
\label{fig:pipeline}
\end{figure}

\subsection{Project Knowledge Discovery}\label{sec:knowledge}

Phase~1 builds the project-specific vocabulary the downstream
linkers need.
Its two analyses are mutually independent, so \approach runs them
concurrently and emits two outputs: an \emph{ambiguity map} over the
component set and an \emph{alias table} over the document.
The alias table is fed into every framing of the \multiframing and
into \coref; the ambiguity map gates the unified validation's
generic-word precheck.

\subsubsection{Model Understanding}\label{sec:model-understanding}

The model-understanding analyser asks a single question of the
component list: which component names can be confused with ordinary
technical English?
A sentence containing the word ``storage,'' for example, does not
necessarily refer to a component named \texttt{Storage}; it may
describe a general storage concern.
We submit the component list to a few-shot LLM classifier that splits
the names into \emph{architectural} names (unambiguous identifiers
such as \texttt{TaskScheduler}) and \emph{ambiguous} names
(single-word identifiers that double as general vocabulary,
e.g., \texttt{Core}, \texttt{Logic}).
A structural filter keeps only single-word, lowercaseable names in
the ambiguous bucket, so that CamelCase compounds, multi-word names,
and acronyms are never classified as generic merely because their
sub-words are common.
The ambiguity map is consumed by the unified validation's
generic-word precheck, thereby giving every multi-framing candidate
the same view of which names need careful handling.

\subsubsection{Document Understanding}\label{sec:doc-understanding}

The document-understanding analyser learns the project-specific
vocabulary that the document uses to refer to components.
It builds an \emph{alias table} whose entries map an observed term
to a canonical component and label the term as either \emph{global}
(safe to recognise anywhere in the document, e.g., a parenthetical
abbreviation or a CamelCase identifier) or \emph{local} (a single
lower-case word whose referent depends on the surrounding context).

\textbf{Alias extraction.}
A first LLM pass scans the document and proposes two kinds of
mappings: \emph{abbreviations} (short forms introduced by the text,
e.g., ``Abstract Syntax Tree (AST)'') and \emph{synonyms}
(alternative names that identify exactly one component, including
trailing words of multi-word component names when the document uses
the trailing word alone).
The prompt asks for an alias scope; the scope is load-bearing
because only the global subset is injected into the three framings
and into \coref, thereby preventing a generic local word from being
broadcast into every sentence as if it were a project-wide name.

\textbf{Alias judging.}
A second LLM pass judges the proposed mappings with an
approval-biased rubric.
The rubric approves abbreviations formed from a component's
initials, unshared trailing words of multi-word names, CamelCase
identifiers, and multi-word phrases containing the component name;
it rejects only terms that are clearly generic, clearly name a
different component, clearly refer to the whole system, or name an
architectural tier or technology platform that encompasses
multiple elements.
The asymmetry is deliberate: a falsely rejected alias permanently
removes a possible reference form, while a falsely approved alias
can still be filtered by the unified validation.

\subsection{The \multiframing}\label{sec:multiframing}

The first observation behind \approach is that, by the definition
of a trace link, the linked component must be present in the linked
sentence in some surface form.
The same link, however, manifests under three complementary
linguistic forms --- an explicit mention of the component name, an
actor-role reference to the component as the grammatical agent of
an action, and an entity-alias reference through a project-specific
synonym or partial name --- and a single LLM prompt biased toward
any one of those forms misses references in the other two.
The \multiframing runs three framings that each target one of the
three reference forms, unions their outputs to recover the
recall that no single framing achieves on its own, and validates
the union through a single evidence-bundle pass that is the sole
quality gate for the entire multi-framing pool.

The three framings are parallel, alias-aware, and intentionally
redundant.
Parallelism keeps the latency of the multi-framing equal to the
slowest framing rather than to their sum.
Alias-awareness gives every framing access to the document alias
table from \S\ref{sec:doc-understanding}, so that a reference such
as ``the dispatcher'' is recognised regardless of which framing
the linker is currently asking about.
Redundancy is the safety margin: if one framing misses a true
reference, another framing recovers it; if one framing hallucinates
a candidate, the unified validation rejects it.

\subsubsection{\linkerA: Explicit-Mention Framing}\label{sec:linkerA}

\linkerA is the mention-framed framing.
A single batched LLM pass asks, for every sentence in the document,
which components are explicitly mentioned or referenced by name or
by a known alias.
The prompt is biased toward inclusion: the framing is intended to
fire whenever the surface form of the component appears, even when
the component is not the grammatical subject of the sentence.
\linkerA therefore covers references such as ``configures the
\texttt{EventDispatcher}'' or ``the \texttt{EventDispatcher}'s
queue is bounded,'' which an actor-framed prompt would skip.

\subsubsection{\linkerB: Actor-Role Framing}\label{sec:linkerB}

\linkerB is the actor-framed framing.
A separate batched LLM pass asks, for every sentence, which
components are architecturally relevant to the sentence --- that
is, whose role, behaviour, interaction, or responsibility is
being described.
The prompt is biased toward agency: it fires when the component
\emph{does} something, \emph{provides} something, or \emph{takes
part in} an interaction.
\linkerB therefore covers references such as ``the
\texttt{EventDispatcher} routes events to subscribers,'' where the
component is the grammatical subject, and the
``handled by \texttt{EventDispatcher}'' / ``via the dispatcher''
prepositional mentions that lexical matchers find awkward.

\subsubsection{\linkerC: Entity-Alias Framing}\label{sec:linkerC}

\linkerC is the entity-framed framing.
Two independent LLM passes scan the document with the global alias
table and ask for every component reference --- by name, by alias,
by partial name, or as a participant in a described interaction.
\approach keeps only the intersection of the two passes; using the
same prompt twice avoids mixing prompt-specific biases, so the
residual disagreement reflects only LLM non-determinism.
\linkerC therefore covers references that surface only through an
alias --- ``the dispatcher,'' ``the AST,'' ``the routing layer'' ---
and that neither \linkerA nor \linkerB picks up because the
canonical name itself is absent.

\subsubsection{Union Merge and Unified Validation}\label{sec:unified}

\textbf{Union merge.}
The three framings emit three candidate sets keyed by
sentence--component pairs.
\approach takes their set union: a pair is admitted to the
candidate pool if any framing emitted it.
Union is deliberate.
Intersection would discard a true reference whenever it surfaces
under only one framing --- e.g., an alias-only reference seen by
\linkerC but not by \linkerA or \linkerB --- and that loss
empirically dominates the false-positive cost of union when a
strong downstream validation is available.\todo{cite RQ: union
vs k=2 voting (s17b $-$4.1pp vs union; s17c with union $\approx$
s15)}

\textbf{Unified evidence-bundle validation.}
Each union candidate carries an \emph{evidence bundle} recording
the matched span, the mention type (proper-case standalone,
lower-case, dotted-path occurrence, alias, or indirect), the
preceding sentence, up to five anchor sentences with confirmed
standalone mentions of the component elsewhere in the document,
the ambiguity flag for the component, and the framing that
proposed the candidate.
Validation therefore sees not only a sentence and a component name
but also \emph{why} the proposer believed the link was plausible,
and which framing's bias it carries.

\textbf{Generic-word precheck.}
For candidates whose component is in the ambiguity map and whose
mention appears only in lower-case form, \approach first asks
whether the word is being used as a component reference or as
ordinary English, conditioned on the up-to-five anchor sentences.
Only candidates classified as component references continue to the
two-pass step.
If the precheck cannot parse a clear answer, \approach keeps the
candidates, thereby preserving the recall-first behaviour used
throughout the pipeline.

\textbf{Two-pass approve--reject.}
Remaining candidates are judged by two independent validation
passes with complementary foci.
The first pass checks \emph{architectural participation}: does the
sentence name the component as a participant that performs an
operation, provides a service, is configured, or takes part in
system behaviour?
The second pass checks \emph{referential specificity}: does the
mention identify this specific architectural element rather than a
generic technical term?
A candidate is accepted only when both passes approve it, thereby
making the unified validation the sole quality gate for every
multi-framing link emitted by \approach.

\subsection{\coref: The Pronominal Linker}\label{sec:coref}

A trace link whose evidence is a discourse pronoun is, by
construction, invisible to the \multiframing: no surface name and
no alias appears in the linked sentence.
\coref targets exactly that gap.
It is supplementary because the gap itself is small in absolute
terms, but the pronominal residue contains the architecturally
densest sentences --- the sentences that describe what a previously
introduced component \emph{does} --- and discarding them costs
real recall.

\textbf{Cases-in-context extraction.}
\coref scans the document for sentences containing a pronoun
(\emph{it}, \emph{they}, \emph{this}, \emph{these}, \emph{that},
\emph{those}, \emph{its}, \emph{their}) or a role-referential
noun phrase (``the dispatcher,'' ``the service,'' ``the module'')
formed from a component-specific terminal word.
A separate LLM call classifies, once per project, which terminal
words of multi-word component names are specific enough to serve
as unambiguous role references; only those terminals are admitted
to the role-referential pattern.
Each candidate sentence is then presented to the resolver with a
$\pm 5$-sentence context window, with the target sentence marked
explicitly, and the resolver is asked to return only the
resolutions it is certain about, citing the target sentence, the
pronoun or role-referential phrase, the component, the antecedent
sentence number, and the antecedent quote.

\textbf{Cited-antecedent check.}
\approach verifies every cited antecedent deterministically before
the resolution reaches validation.
For canonical-name antecedents, the antecedent sentence must
contain a standalone mention of the component name.
For alias-based antecedents, the resolver must explicitly mark the
antecedent as alias-based, in which case the alias-table lookup
provides the same deterministic check.
Citation grounding therefore removes any antecedent the LLM may
have fabricated, thereby reducing pronoun resolution from an
open-ended reasoning task to a deterministic verification.

\textbf{Coref-specific validation.}
Surviving resolutions then pass through a single LLM validation
call whose focus is coref-specific: does the pronoun or
role-referential phrase in this sentence actually refer to the
named component as an architectural participant --- performing
operations, providing services, or being the grammatical topic of
the sentence?
This validation is the s17e change: in earlier variants, coref
links bypassed every quality gate and entered the final link set
on the strength of the antecedent check alone, which empirically
let through false positives whose pronouns resolved to a component
by surface association rather than by architectural role.
The added pass closes that hole and is, on the v2.6.2 benchmark,
responsible for a $+3.2$~pp F1 gain and a drop in false-positive
coref links from 43 to 14.\todo{cite RQ: s17e vs s15 ablation
showing the +3.2pp / FP 43->14 attribution.}

\subsection{Link Consolidation}\label{sec:merge}

The final phase is intentionally simple.
The \multiframing has emitted a validated link set through its
unified evidence-bundle validation, and \coref has emitted a
validated link set through its cited-antecedent check plus its
coref-specific validation.
\approach concatenates the two sets and removes duplicate
sentence--component pairs in a fixed priority order: multi-framing
first, coref second.
The priority preserves the most direct evidence whenever both
linkers recover the same pair and reflects the linker hierarchy ---
the \multiframing is primary because its three framings exhaust
the surface forms of textual reference admitted by the trace-link
definition, while \coref contributes the pairs whose evidence is
a pronoun and which the multi-framing pipeline, by the same
definition, cannot see.

The final output is a documentation-to-model trace-link set; these
links can be evaluated directly at the architecture level or
composed with existing model-to-code mappings to recover
documentation-to-code traces (\S\ref{sec:experiment}).
